\documentclass[master]{thesis-uestc}

\title{中文题目}{English Title} % 论文题目
\author{姓名}{English Name} % 作者姓名
\setdate[submit]{2022年3月17日} % 论文提交日期，可留空
\setdate[oral]{2022年4月15日}  % 答辩日期，可留空
\setdate[confer]{2022年6月} % 学位授予日期，可留空
\advisor{导师姓名\chinesespace 导师职称}{English name English title}
% \coAdvisor{合作导师姓名\chinesespace 导师职称}{Co advisor English name English title} % 仅专业硕士/博士使用，在扉页/英文首页添加合作导师，不使用请注释
\school{计算机科学与工程学院(网络空间安全学院)}{School of Computer Science and Engineering(School of Cyberspace Security)} % 学院信息
\major{计算机科学与技术}{Computer Science and Technology} % 专业信息
\studentnumber{xxxxx} % 学号
\ProfessionalDegreeArea{随便学学} % 专业硕士专用：专业学位领域
\ClassificationNumber{TP309.2} % 分类号
\ClassifiedClass{公开} % 密级
\UDCNumber{004.78} % UDC号
\Chairman{xxxxx} % 答辩委员会主席

% 取消注释以下内容，用于禁止文中换行处的英语单词自动截断换行。
% \tolerance=1
% \emergencystretch=\maxdimen
% \hyphenpenalty=10000
% \hbadness=10000

\makeglossaries % 产生缩略词表/符号表专用,不使用时请注释
\newacronym[description=逻辑卷管理器]{lvm}{LVM}{Logical Volume Manager} % 定义缩略词:以本项为例,逻辑卷管理器为中文名称;lvm用于文内引用;LVM为显示的应为缩略语或符号;Logical Volume Manager为显示的英文全称/描述。
\newglossaryentry{tree}{name={tree}, description={trees are the better humans}}  % 定义符号:以本项为例,name={tree}为符号名称;tree用于文内引用; description={trees are the better humans}为显示的描述。页码自动添加。

\begin{document}
\makecover % 封面+中英文扉页
\originalitydeclaration % 原创新声明
% \signatureofdeclaration{signature.pdf} % 用于添加扫描版签字后的原创新声明（使用时取消注释本行，并注释掉上一行）
% 中文摘要
\begin{chineseabstract}

    \chinesekeyword{xxx，xxx，xxx} % 中文关键词
\end{chineseabstract}
% 英文摘要
\begin{englishabstract}

    \englishkeyword{xxx, xxx, xxx} % 英文关键词
\end{englishabstract}

\thesistableofcontents % 目录
\thesisfigurelist % 图目录，仅在需要时添加，一般情况下请注释
\thesistablelist % 表目录，仅在需要时添加，一般情况下请注释
% \glsaddall % 默认仅显示被正文引用的项,取消注释以显示所有已定义的缩略词/符号
\thesisglossarylist % 缩略词表，仅在需要时添加，一般情况下请注释
\thesissymbollist % 符号表，仅在需要时添加，一般情况下请注释

% 正文内容
\chapter{绪\hspace{6pt}论}

角标参考文献\citing{chen2001hao}，普通参考文献\cite{clerc2010discrete}。

这是符号\gls{tree}。

这是缩略词\acrlong{lvm}的长引用，这是缩略词的短引用\acrshort{lvm}。

\begin{figure}
    \includegraphics[width=0.5\linewidth]{pic/logo.pdf}
    \caption[short catption 1]{Test caption 1}
\end{figure}


\chapter{相关研究基础}
\chapter{工作内容xxx}
\chapter{全文总结与展望}

算法框：

\begin{algorithm}[H]
    \KwData{this text}
    \KwResult{how to write algorithm with \LaTeX2e}
    initialization\;
    \While{not at end of this document}{
        read current\;
        \eIf{understand}{
            go to next section\;
            current section becomes this one\;
        }{
            go back to the beginning of current section\;
        }
    }
    \caption{How to wirte an algorithm.}
\end{algorithm}

\thesisacknowledgement

xxxx % 直接填写致谢内容，写法与正文一致

\thesisbibliography[large]{reference} % 参考文献

\thesisappendix
\chapter{xxxx} % 直接填写附录内容，写法与正文一致

% 攻读学位期间成果（本科不添加），例如：
\begin{thesistheaccomplish}
    \section{学术论文}
    \bibitem{SGXDedup} \textbf{Ren, Yanjing} and Li, Jingwei and Yang, Zuoru and Lee, Patrick PC and Zhang, Xiaosong. Accelerating Encrypted Deduplication via SGX[C]. Proc.of USENIX ATC, 2021, 957-971. \textbf{CCF-A}
    \section{发明专利}
    \bibitem{CN111338572B} 李经纬, 杨祚儒, \textbf{任彦璟}, 李柏晴, 张小松. 一种可调节加密重复数据删除方法:CN111338572B[P]. 2021-09-14.
\end{thesistheaccomplish}

% 以下为本科同学专用，插入文献翻译
\thesistranslationoriginal
\section{Tahoe-LAFS: The Least-Authority File System}
% \insertPDFPage{} % 用于插入单页PDF文件，例如原始文献（该操作会导致页码被取消，请谨慎使用）

\thesistranslationchinese
\section{Tahoe-LAFS：最小权限文件系统}

\end{document}